\mxcorrections{4}\label{C04-CORR}
\textbf{Usage professeur.} Les réponses qui suivent concernent les tâches non résolues dans le cours. Faire d’abord expliciter la démarche ; fournir une aide graduée avant la solution. Les durées sont des estimations à étalonner en classe.
\subsection*{Diagnostic et remédiation}
\begin{enumerate}
\item $\dfrac{8-2}{3-1}=3$. \textit{Action :} Représenter les accroissements horizontaux et verticaux.
\item $4h+h^2$. \textit{Action :} Développer avant de simplifier.
\item $4+h$. \textit{Action :} Factoriser h au numérateur avant la division.
\item $0$ et $0{,}1$. \textit{Action :} Vérifier par élévation au carré.
\item $3$ et $2$. \textit{Action :} Représenter la distance à 0.
\item $y=2x+3$. \textit{Action :} Substituer les coordonnées pour trouver l’ordonnée à l’origine.
\end{enumerate}
\subsection*{Activité de recherche}
Les pentes sont $2+h$ ; limite 2. La droite passe par $(1;1)$, donc $y=2x-1$. Le dessin motive, le calcul justifie.
\subsection*{Applications immédiates}
\begin{enumerate}
\item $f(2+h)-f(2)=7h+h^2$ ; taux $7+h$ pour $h\ne0$.
\item Le taux vaut $4+3h$, donc le nombre dérivé vaut 4.
\item $y=5-2(x-3)=-2x+11$ ; approximation $5-2(0{,}04)=4{,}92$.
\item Le taux en 2 est $\dfrac{|h|}{h}$, avec les deux valeurs limites opposées. Pour la racine en 0, le taux $\dfrac{1}{\sqrt h}$ pour $h>0$ dépasse toute borne lorsque h se rapproche de 0 : aucun nombre dérivé réel.
\end{enumerate}
\subsection*{Exercices : démarche et résultats}
\paragraph{Exercice 4.1 -- Deux sécantes.} Pentes 3 et 2{,}1 ; nombre dérivé 2. La seconde est plus proche de la pente tangente ; il s’agit ici d’un calcul exact de taux.
\paragraph{Exercice 4.2 -- Une tangente.} Taux $4a+2h=-4+2h$, donc $f\prime(-1)=-4$. $f(-1)=3$ ; tangente $y=3-4(x+1)=-4x-1$.
\paragraph{Exercice 4.3 -- Coût marginal.} $C(10)=400$, $C(11)=441$, accroissement 41 pour une unité. Le taux autour de 10 vaut $40+h$, donc dérivée 40 : approximation de l’accroissement pour une variation petite, pas égalité pour une unité entière.
\paragraph{Exercice 4.4 -- Lire une tangente.} Sur la figure du cours : points (1;0) et (2;4), pente 4. Dans la nouvelle situation, $f(2)=3$, pente $-\dfrac{1}{2}$. Approximation $3-0{,}5\times0{,}1=2{,}95$. L’ordonnée de B n’est pas nécessairement $f(4)$.
\paragraph{Exercice 4.5 -- Contrôler une approximation.} $1+2(-0{,}02)=0{,}96$. Exact $0{,}9604$ ; exact moins approché $0{,}0004>0$. Le terme omis est $h^2$ : approximation par défaut.
\paragraph{Exercice 4.6 -- Tangente et intersection.} La tangente a pour équation $y=1+3(x-1)=3x-2$. Le développement du produit donne bien $x^3-3x+2$. Les abscisses des intersections sont donc 1 et $-2$, par la règle du produit nul. Les points sont $(1;1)$ et $(-2;-8)$. Il existe deux points communs : cela suffit à réfuter l’affirmation générale.
\paragraph{Exercice 4.7 -- Sécantes numériques.} Liste des taux $2+h$ : 2{,}1 ; 2{,}01 ; 2{,}001 et 1{,}9 ; 1{,}99 ; 1{,}999. Le calcul algébrique donne le réel 2. À très petit h, le calcul flottant peut perdre de la précision.
\paragraph{Exercice 4.8 -- Reprise différée.} Pour $h\ne0,-1$, le taux est $-\dfrac{1}{1+h}$ ; il tend vers $-1$. Approximation $1-0{,}02=0{,}98$. La valeur exacte $\dfrac{50}{51}$ permet un contrôle.
\paragraph{CQFD, exercice 17 p. 127, questions toutes.} Les abscisses de A, B et C sont 0, 1 et 2. Les tangentes donnent respectivement les pentes 2, 0 et $-2$. Ainsi $f\prime(0)=2$, $f\prime(1)=0$, $f\prime(2)=-2$. Les unités des axes doivent être lues sur le quadrillage.
\paragraph{CQFD, exercice 19 p. 127, questions 1 et 2.} Le point A se lit $(1;4)$, donc $f(1)=4$. La tangente passe par $(0;1)$ et $(1;4)$ : pente 3. Ainsi $f\prime(1)=3$ et la tangente a pour équation $y=4+3(x-1)=3x+1$.
\subsection*{QCM : réponses et raisons}
\textbf{QCM 1.}
\begin{enumerate}
\item B. $\dfrac{4-1}{2-1}=3$.
\item A. $4+2(x-1)$.
\item D. Il s’agit d’une valeur en un point.
\item C. Taux $6+h$.
\item B. Fonction affine tangente.
\end{enumerate}
\textbf{QCM 2.}
\begin{enumerate}
\item C. Taux $\dfrac{1}{\sqrt h}$ non fini.
\item A. $|h|=-h$.
\item D. Ordonnée constante.
\item B. Une variation de position divisée par une variation de temps s’exprime en $\mathrm{m\,s^{-1}}$.
\item C. Égalité exacte pour une affine, approximation en général.
\end{enumerate}
\subsection*{Vrai ou faux}
\begin{enumerate}
\item Faux : $x^3$ et la tangente $3x-2$ en 1 se rencontrent aussi en $-2$.
\item Faux : les taux de part et d’autre valent $-1$ et 1.
\item Vrai : le taux vaut exactement m pour tout h non nul.
\end{enumerate}
\paragraph{Différenciation.} Pour un blocage technique, proposer la première étape seulement (domaine, facteur ou univers de référence), puis retirer l’aide. Pour un élève autonome, demander une seconde méthode et le traitement d’un cas limite. À l’oral, autoriser une reformulation initiale dans la langue maîtrisée, puis exiger une phrase mathématique précise en français.
